% Source: https://artofproblemsolving.com/community/c3h2528320_function
% Stable identity: aops-topic-2528320-q01
\textbf{(I)} For $a=-\dfrac34$,
\[f'(x)=-\dfrac{3}{4x}+\dfrac{1}{2\sqrt{1+x}}=\dfrac{2x-3\sqrt{1+x}}{4x\sqrt{1+x}}.\]
For $x>0$: $2x>3\sqrt{1+x}\iff4x^2>9+9x\iff(x-3)(4x+3)>0\iff x>3$. So $f$ is \textbf{decreasing on $(0,3)$} and \textbf{increasing on $(3,+\infty)$}.

\textbf{(II)} \emph{Necessity.} Take $x=1$: $\sqrt2\le\dfrac{1}{2a}$, which forces $a>0$ and $a\le\dfrac{1}{2\sqrt2}=\dfrac{\sqrt2}{4}$.

\emph{Sufficiency.} Let $0<a\le\dfrac{\sqrt2}{4}$ and $t=\dfrac1a\ge2\sqrt2$. Multiplying the inequality by $t>0$, it becomes $\ln x+t\sqrt{1+x}\le\dfrac{t^2\sqrt x}{2}$, i.e.
\[g(t)=\dfrac{\sqrt x}{2}\,t^2-\sqrt{1+x}\,t-\ln x\ge0.\]
For fixed $x$, $g$ is a quadratic in $t$ opening upwards with vertex at $t^*=\sqrt{\dfrac{1+x}{x}}=\sqrt{1+\dfrac1x}$.

\emph{Case $x\ge\dfrac17$.} Then $t^*\le\sqrt8=2\sqrt2\le t$, so $g(t)\ge g\left(2\sqrt2\right)=4\sqrt x-2\sqrt2\sqrt{1+x}-\ln x=:p(x)$. We show $p(x)\ge0$ for all $x>0$:
\[p'(x)=\dfrac{2}{\sqrt x}-\dfrac{\sqrt2}{\sqrt{1+x}}-\dfrac1x=\dfrac{N}{x\sqrt{1+x}},\qquad N=\sqrt{1+u^2}\,(2u-1)-\sqrt2\,u^2,\quad u=\sqrt x.\]
If $u\le\dfrac12$, then $N<0$. If $u>\dfrac12$, both terms of $N$ are positive, and $N$ has the sign of
\[\left(1+u^2\right)(2u-1)^2-2u^4=(u-1)\left(2u^3-2u^2+3u-1\right).\]
The cubic $2u^3-2u^2+3u-1$ is increasing (its derivative $6u^2-4u+3$ has negative discriminant) and equals $\dfrac14$ at $u=\dfrac12$, so it is positive for $u>\dfrac12$. Hence $p'<0$ on $(0,1)$ and $p'>0$ on $(1,\infty)$, and $p(x)\ge p(1)=4-4-0=0$.

\emph{Case $\dfrac{1}{e^2}\le x<\dfrac17$.} Now $t^*>2\sqrt2$, and $g(t)\ge g(t^*)=-\ln x-\dfrac{1+x}{2\sqrt x}=:s(x)$. Then
\[s'(x)=-\dfrac1x-\dfrac{x-1}{4x\sqrt x}=-\dfrac{x+4\sqrt x-1}{4x\sqrt x}.\]
As $e^2<9$, $x\ge\dfrac1{e^2}>\dfrac19$, so $\sqrt x>\dfrac13$ and $x+4\sqrt x-1>\dfrac43-1>0$. So $s$ is decreasing and $s(x)\ge s\left(\dfrac17\right)=\ln7-\dfrac{4}{\sqrt7}$. Finally $\ln7>\dfrac74$ (because $e^7<3^7=2187<2401=7^4$) and $\dfrac{4}{\sqrt7}<\dfrac74$ (because $16^2=256<343=49\cdot7$), so $s(x)>0$.

In both cases $g(t)\ge0$, so the inequality holds for all $x\ge\dfrac{1}{e^2}$. Therefore the range is
\[\boxed{0<a\le\dfrac{\sqrt2}{4}}.\]
